\documentclass[10pt,a4paper]{amsart}
\usepackage[margin=19mm]{geometry}
\usepackage{amsmath,amssymb,mathtools}
\usepackage[scr=rsfs,cal=euler]{mathalfa}
\usepackage{xfrac}
\usepackage[unicode,hidelinks,bookmarks=false]{hyperref}
\newcommand{\ie}{i.e.\ }
\newcommand{\cf}{cf.\ }
\newcommand{\resp}{respectively\ }
\newcommand{\Subsection}{\subsection}
\providecommand{\nobreakdash}{\nobreak\hbox{-}}
\setlength{\emergencystretch}{2em}
\multlinegap0pt
\usepackage{bm}
\usepackage[shortlabels]{enumitem}
\usepackage{cancel}
\usepackage{mathtools}
\usepackage{xcolor}
\usepackage{float}
\usepackage{amssymb}
\usepackage[normalem]{ulem}
\usepackage{tensor}
\usepackage{tikz}
\usepackage{csquotes}
\newtheorem{lemma}{Lemma}
\newtheorem{theorem}{Theorem}
\newcommand{\Z}{{\mathbb Z}}
\newcommand{\be}{\begin{equation}}
\newcommand{\ee}{\end{equation}}
\newcommand{\id}{\text{id}}
\newcommand{\mc}[1]{\mathcal{#1}}
\newcommand{\paren}[1]{\left( #1 \right)}
\title{Non-Invertible Symmetries Mixing with Witt Non-Trivial \\ Quantum Cellular Automata}
\author{Kansei Inamura, Oskar Wojdel, Lukasz Fidkowski, Sakura Schafer-Nameki}
\date{Source: arXiv:2607.21698v2 (CC BY 4.0)}
\begin{document}
\maketitle

\noindent\emph{Source and licence.} Non-Invertible Symmetries Mixing with Witt Non-Trivial \\ Quantum Cellular Automata. Version arXiv:2607.21698v2 (CC BY 4.0), distributed by arXiv under the Creative Commons Attribution 4.0 International licence, http://creativecommons.org/licenses/by/4.0/, as recorded in the arXiv licence metadata for this exact version and shown on the abstract page https://arxiv.org/abs/2607.21698v2. Full licence notice: \texttt{licenses/arxiv-2607.21698-CC-BY-4.0.txt}.\par\smallskip

\noindent\textit{Editorial scope.} Counted unit: the article's theorem thm:QCA\_extendability (source0.tex lines 7032--7034). The article's complete original proof of it is delivered with it (source0.tex lines 7036--7188, one \texttt{proof} environment of 153 lines). No mathematics is written, repaired or re-derived: the statement and the proof are the article's own lines, in the article's own notation, and no retained line is substituted or re-flowed.  Retained uncounted as the source-local context the proof needs: lines 6817--6822.  Nothing was omitted from any retained span, so the package carries no omission record.  No retained line contains a citation, so the wrapper carries no bibliography and no bibliography entry.  No retained line contains a figure environment, an image-inclusion command or drawing code, so no image is delivered and the package declares no asset dependency.  Editorial formatting only, in addition: an AMS \texttt{amsart} preamble that restates the article's own declarations and macros used by the retained lines, namely the environments \texttt{lemma}, \texttt{theorem} on separate counters, together with the standard packages those lines need.  The retained environments print in this document as Lemma 1, Theorem 1 in order of appearance, whereas the article prints the same items with the numbering of its own full document.  The complete native source of the article as distributed in the arXiv e-print for this exact version is bundled unchanged as \texttt{sources/205887/source0.tex}.
\begin{lemma}\label{lemma:exact_support}
 For any $a + \mathfrak{I} \in \mc{A}_{\Z_2^{(1)}} / \mathfrak{I}$, $\exists I \in \mathfrak{I}$ such that
 \be
 \operatorname{Supp}(a + I) = \operatorname{Supp}(a + \mathfrak{I})
 \ee
\end{lemma}

\begin{theorem}\label{thm:QCA_extendability}
 Let $\alpha : \mc{A}_{\Z_2^{(1)}} / \mathfrak{I} \rightarrow \mc{A}_{\Z_2^{(1)}} / \mathfrak{I}$ be a QCA such that $\alpha(Z_{\delta \bm{e}} + \mathfrak{I}) = Z_{\delta \bm{e}} + \mathfrak{I}$. Then, there exists an extension $\tilde{\alpha}$ to the full algebra. Furthermore, we may choose it such that $\tilde{\alpha}(Z_f) = Z_f \; \forall f \in F$.
\end{theorem}

\begin{proof}
 The proof is constructive. Let us define $\tilde{\alpha}(Z_f) = Z_f$.

 Now choose $x_f \in \alpha(X_f + \mathfrak{I})$ arbitrarily. We will later see how to make the choice correctly. Note that we have:
 \begin{align}
 (x_f + \mathfrak{I}) (Z_{\delta \bm{e}} + \mathfrak{I}) &= \alpha(X_f Z_{\delta \bm{e}} + \mathfrak{I}) \nonumber \\
 (Z_{\delta \bm{e}} + \mathfrak{I}) (x_f + \mathfrak{I}) &= \alpha(Z_{\delta \bm{e}} X_f + \mathfrak{I})
 \end{align}
 so the commutation relations of $x_f$ and $Z_{\delta \bm{e}}$ are the same as those of $X_f$ and $Z_{\delta \bm{e}}$ (or rather, their equivalence classes). This means that $X_f x_f$ commutes with all $Z_{\delta \bm{e}}$, or in other words
 \be
 X_f x_f Z_{\delta \bm{e}} - Z_{\delta \bm{e}} X_f x_f = J_{e, f} \in \mathfrak{I} \qquad \forall e, f \,.
 \ee

 From this we can also see that $Z_{\delta \bm{e}} J_{e, f} Z_{\delta \bm{e}} = - J_{e, f}$, and therefore
 \be
 \paren{X_f x_f + \frac{1}{2} Z_{\delta \bm{e}} J_{e, f}} Z_{\delta \bm{e}} = Z_{\delta \bm{e}} \paren{X_f x_f + \frac{1}{2} Z_{\delta \bm{e}} J_{e, f}} \,.
 \ee

 That is, for every edge $e \in E$, there is some $I \in \mathfrak{I}$, such that $X_f x_f + I$ commutes with $Z_{\delta \bm{e}}$. This means that $e \notin \operatorname{Supp}(X_f x_f + \mathfrak{I})$, and so
 \be
 \operatorname{Supp}(X_f x_f + \mathfrak{I}) \subseteq F \,.
 \ee

 It can be shown (Lemma \ref{lemma:exact_support}) that there is always some $a_f \in X_f x_f + \mathfrak{I}$ whose support is exactly the same as that of its equivalence class:
 \be
 \operatorname{Supp}(a_f) = \operatorname{Supp}(X_f x_f + \mathfrak{I}) \subseteq F \,.
 \ee
 In fact, we may choose a representative $a_f$ which is written entirely in terms of $Z_{\delta \bm{e}}$. Any $X_f$, would have to appear in configurations without boundary, but those are just symmetry operators, which can all be eliminated by choosing a different representative:
 \be
 \prod_{f \in c_2} X_f \prod_{e \in c_1} Z_{\delta \bm{e}} = \prod_{e \in c_1} Z_{\delta \bm{e}} - \paren{1 - \prod_{f \in c_2} X_f} \prod_{e \in c_1} Z_{\delta \bm{e}} \,,
 \ee
 where due to $\partial c_2 = 0$, the second term is in $\mathfrak{I}$.

 Finally, we define $\tilde{\alpha}(Z_f) = Z_f$ and $\tilde{\alpha}(X_f) = X_f a_f$.

 Let us check that such $\tilde{\alpha}$ really is a $*$-homomorphism. Clearly $\tilde{\alpha}(Z_f) = Z_f$ all commute between each other. Furthermore, $\tilde{\alpha}(Z_f), \tilde{\alpha}(X_f)$ commute or anticommute as they should: the sign comes from the $X_f$ in $\tilde{\alpha}(X_f) = X_f a_f$, and $[a_f, Z_p] = 0$ because $a_f$ is written entirely in terms of $Z$.

 Now, we must verify that $\tilde{\alpha}(X_f)$ commute between each other. Write
 \be
 a_f = \sum_{b^2} A^f_{b^2} \prod_{f \in b^2} Z_f \,.
 \ee

 Then
 \be
 X_f a_f X_p a_p = X_f X_p \paren{\sum_{b^2} (-1)^{\bm{b^2}(p)} A^f_{b^2} \prod_{f \in b^2} Z_f } a_p \,.
 \ee

 Similarly
 \be
 X_p a_p X_f a_f = X_p X_f a_f \paren{ \sum_{b'^2} (-1)^{\bm{b}'^2(f)} A^p_{b'^2} \prod_{f \in b'^2} Z_f } \,.
 \ee

 Due to the fact that $\alpha(X_f + \mathfrak{I})$ and $\alpha(X_p + \mathfrak{I})$ commute, we also must have $X_f a_f X_p a_p - X_p a_p X_f a_f \in \mathfrak{I}$, or equivalently if we multiply by $X_f X_p$ on the left,
 \be
 (X_p a_f X_p) a_p - a_f (X_f a_p X_f) \in \mathfrak{I} \,,
 \ee
 where $X a X$ are just the expressions in parentheses earlier.
 
 \textbf{Claim.} $\langle Z_{\delta \bm{e}} \rangle \cap \mathfrak{I} = \{ 0 \}$.

 We see that that difference, which is in $\mathfrak{I}$, can also be written entirely in terms of $Z_{\delta \bm{e}}$ (conjugating by $X_f$ does not change this, it only adds some signs). By the claim, the difference is $0$, so $\tilde{\alpha}(X_f)$ commute between each other.

 \textbf{Proof.} Clearly, every element of $\langle Z_{\delta \bm{e}} \rangle \subset \mc{A}$ commutes with every $Z_f$.

 We are actually going to show something slightly stronger. That if some $I \in \mathfrak{I}$ also commutes with all $Z_f, \forall f \in F$, then $I = 0$.

 We argue by contradiction: suppose there is some $I \in \mathfrak{I} \setminus \{0\}$ which commutes with all $Z_f$. Without loss of generality, we may assume that $I$ is written entirely in terms of $X$ operators. Indeed, expanding $I$ in the basis of $\mc{A}$, $\prod X \prod Z$ with the products running over arbitrary 2-chains, we can group the terms according to their configuration of $Z$ operators. In such an expression, each ``coefficient'' (which is actually an operator written entirely in terms of $X$ operators) must commute with all $Z_f$. At least one of those coefficients is non-zero, and we may pick that one.

 But then, if we have
 \be
 I = \sum_{c_2} I_{c_2} \prod_{f \in c_2} X_f \,,
 \ee
 this is still written as a linear combination of linearily independent terms. For any $c'_2 \neq 0$, choose a plaquette $p \in c'_2$:
 \be
 I = Z_p I Z_p = - I_{c'_2} \prod_{f \in c'_2} X_f + \sum_{c_2 \neq c'_2} (-1)^{\bm{p}(c_2)} I_{c_2} \prod_{f \in c_2} X_f \,,
 \ee
 so $I_{c'_2} = 0$. The only term which remains is $I = I_0 \id{}$. But this is invertible for any $I_0 \neq 0$, so we conclude $I = 0$. With this, the proof of the claim is complete.
 
 We have shown that
 \begin{align}
  \tilde{\alpha}(Z_f) &= Z_f \,, & \tilde{\alpha}(X_f) &= X_f a_f \,,
 \end{align}
 preserves all the commutation relations of the generators. $Z_f$ is also self-adjoint, and unitary (i.e. squares to 1).

 To show that $X_f a_f$ is self-adjoint, note that
 \be
 (X_f a_f)^\dag + \mathfrak{I} = \alpha\paren{(X_f + \mathfrak{I})^\dag} = \alpha\paren{X_f + \mathfrak{I}} = X_f a_f + \mathfrak{I} \,.
 \ee

 We manipulate the condition by multiplying by a unitary:
 \be
 (X_f a_f)^\dag - X_f a_f \in \mathfrak{I} \iff a_f^\dag - X_f a_f X_f \in \mathfrak{I} \,.
 \ee

 The operator on the right can clearly be written using only $Z_{\delta \bm{e}}$, so by using the claim it must be $0$. The self-adjointness of $X_f a_f$ holds exactly, not only at the level of equivalence classes.

 Finally, to show that $X_f a_f$ is unitary:
 \be
 (X_f a_f)^2 + \mathfrak{I} = \alpha\paren{(X_f + \mathfrak{I})^2} = 1 + \mathfrak{I} \,,
 \ee
 we manipulate the condition
 \be
 X_f a_f X_f a_f - 1 \in \mathfrak{I} \iff a_f^\dag a_f - 1 \in \mathfrak{I} \,.
 \ee

 By using the claim again, we conclude that this equality must hold exactly.

 Hence, $\tilde{\alpha}$ defines a $*$-homomorphism.

 We can construct its inverse explicitly,
 \begin{align}
  \tilde{\alpha}^{-1}(Z_f) &= Z_f \,, & \tilde{\alpha}^{-1}(X_f) &= X_f a_f^\dag \,,
 \end{align}
 to show that it is a $*$-automorphism.

 The commutative diagram in the statement of the theorem is equivalent to $\pi \circ \tilde{\alpha} = \alpha \circ \pi$, when restricted to $\mc{A}_{\Z_2^{(1)}}$. As a preliminary step, note that
 \be
 X_f \in \langle X_f a_f, Z_{\delta \bm{e}} \; \mid \; f \in F, e \in E \rangle
 \ee
 because $a_f \in \langle Z_{\delta \bm{e}} \rangle$. Hence
 \be
 \tilde{\alpha} \paren{\mc{A}_{\Z_2^{(1)}}} = \mc{A}_{\Z_2^{(1)}} \,,
 \ee
 and so the equality $\pi \circ \tilde{\alpha} = \alpha \circ \pi$ makes sense.
 
 As always we can check equality on the generators:
 \begin{align}
  \alpha(Z_{\delta \bm{e}} + \mathfrak{I}) &= Z_{\delta \bm{e}} + \mathfrak{I} \,, \\
  \tilde{\alpha}(Z_{\delta \bm{e}}) + \mathfrak{I} &= Z_{\delta \bm{e}} + \mathfrak{I} \,,
 \end{align}
 and similarly
 \begin{align}
  \alpha(X_f + \mathfrak{I}) &= x_f + \mathfrak{I} = X_f a_f + \mathfrak{I} \,, \\
  \tilde{\alpha}(X_f) + \mathfrak{I} &= X_f a_f + \mathfrak{I} \,.
 \end{align}

 Lastly, we must show that $\tilde{\alpha}$ is also locality preserving, assuming $\alpha$ has a radius of spread $r$. Obviously we may focus on $X_f a_f$, and all we really need to check is that $a_f$ has bounded support (uniformly bounded for all $f$).

 In the quotient:
 \begin{multline}
  \operatorname{Supp}(a_f + \mathfrak{I}) = \operatorname{Supp}\paren{(X_f + \mathfrak{I}) \alpha(X_f + \mathfrak{I})} \\
  \subseteq f \cup \operatorname{Supp}(\alpha(X_f + \mathfrak{I})) \subseteq B_r(f) \,,
 \end{multline}
 where $B_r(f)$ is the ball of radius $r$ with center $f$. That last inclusion follows from the assumption that $\alpha$ has bounded spread (is locality preserving).

 Now recall that we chose the representative $a_f$ exactly as that one which has the same support as its equivalence class. Hence:
 \be
 \operatorname{Supp}(a_f) = \operatorname{Supp}(a_f + \mathfrak{I}) \subseteq B_r(f) \,,
 \ee
 which also shows that $\tilde{\alpha}$ is locality preserving, and hence an ordinary QCA on a tensor product algebra.

 We remark here that this last step relies on the fact that for $a_f$ in particular, its support in $\mc{A}_{\Z_2^{(1)}}$ as defined in this paper matches the usual notion of support. This is because $a_f$ is written entirely in terms of $Z$ operators, so when checking for where it does not commute with $X_f, Z_{\delta \bm{e}}$, only $X_f$ may fail to commute. This ends up being the same as checking where $a_f$ is not proportional to the identity.
\end{proof}

\end{document}
